\section{Eight-derivative couplings from 12d}
\label{sec:12d_eff_act}

We follow \cite{Minasian:2015bxa} in deriving the 10d type IIB eight-derivative
couplings of gravity and axio-dilaton from a 12d perspective.

%----------------------------------------------------------------------
\subsection{Preliminaries}
%----------------------------------------------------------------------

The two-derivative 10d effective action for the metric $g_{mn}$ and axio-dilaton
$\tau = C_{(0)} + i e^{-\Phi}$ is (in SL(2,$\mathbb{R}$)-invariant Einstein frame)
\begin{equation}\label{eq:S0}
    S^{(0)} = \int \bigl(R - 2\,P_m\,\bar{P}^m\bigr)\,{*_{10}}\,1\,,
\end{equation}
where $R$ is the 10d Ricci scalar and $P_m$ packages the $\tau$ kinetic term
in a way that makes the $\text{SL}(2,\mathbb{R})/\text{U}(1)$ coset structure manifest.

\paragraph{Axio-dilaton quantities.}
Writing $\tau = \tau_1 + i\tau_2$, the complex vielbein, its conjugate,
and the composite U(1) connection are
\begin{equation}\label{eq:P_Q_def}
    P_m = \frac{i}{2\tau_2}\,\nabla_m\tau\,,\qquad
    \bar{P}_m = -\frac{i}{2\tau_2}\,\nabla_m\bar\tau\,,\qquad
    Q_m = -\frac{1}{2\tau_2}\,\nabla_m\tau_1
        = \frac{i}{2}\bigl(P_m - \bar P_m\bigr)\,.
\end{equation}
Covariant derivatives (with respect to both the Levi-Civita and U(1) connections) are
\begin{equation}\label{eq:D_def}
    D_m P_n = \nabla_m P_n - 2i\,Q_m\,P_n\,,\qquad
    D_m \bar P_n = \nabla_m \bar P_n + 2i\,Q_m\,\bar P_n\,.
\end{equation}
The antisymmetrised derivative satisfies
$D_{[m}P_{n]} = 0$ as long as $\tau$ is globally well-defined
(no 7-brane sources).

%----------------------------------------------------------------------
\subsection{12d metric embedding}
%----------------------------------------------------------------------

The key F-theory idea is that the axio-dilaton $\tau$ may be identified
with the complex-structure modulus of an auxiliary torus.
Define a 12d metric $\hat{g}_{MN}$ with $M = 0,\dots,11$ by
\begin{equation}\label{eq:12d_metric}
    \hat{g}_{MN}
    = \begin{pmatrix}
        g_{mn} & 0 & 0 \\[4pt]
        0 & \dfrac{1}{\tau_2} & \dfrac{\tau_1}{\tau_2} \\[6pt]
        0 & \dfrac{\tau_1}{\tau_2} & \dfrac{\tau_1^2+\tau_2^2}{\tau_2}
    \end{pmatrix},
\end{equation}
where $g_{mn}$ is the 10d Einstein-frame metric and the lower-right $2\times2$
block is the flat metric on a torus $T^2$ with modular parameter $\tau$.
The torus has fixed (unphysical) volume, and only 10d derivatives act:
$\partial_M = (\partial_m,\,0,\,0)$.

The two-derivative action \eqref{eq:S0} then lifts to
\begin{equation}\label{eq:S0_12d}
    S^{(0)} = \int \hat{R}\,{*_{10}}\,1\,,
\end{equation}
where $\hat R = \hat R_{MN}\hat g^{MN}$ is the 12d Ricci scalar.
There is no 12d promotion of the 10d measure.

%----------------------------------------------------------------------
\subsection{Non-holomorphic Eisenstein series}
%----------------------------------------------------------------------

At the eight-derivative level, the known parity-even part of the 10d
effective action is proportional to $f_0\,t_8\,t_8\,R^4$.
The SL(2,$\mathbb{Z}$)-invariant, non-holomorphic Eisenstein series of weight $3/2$
that serves as the modular coefficient is
\begin{equation}\label{eq:f0}
    f_0(\tau,\bar\tau)
    = \sum_{(m,n)\neq(0,0)} \frac{\tau_2^{3/2}}{|m + n\tau|^3}\,.
\end{equation}
For large $\tau_2$ (weak coupling, $g_s = e^{\langle\Phi\rangle} \ll 1$) this has the expansion
\begin{equation}\label{eq:f0_expansion}
    f_0(\tau,\bar\tau)
    = 2\zeta(3)\,\tau_2^{3/2}
      + \frac{2\pi^2}{3}\,\tau_2^{-1/2}
      + \mathcal{O}(e^{-\tau_2})\,,
\end{equation}
where the first term is tree-level, the second is one-loop,
and the exponentials are D-instanton contributions.

%----------------------------------------------------------------------
\subsection{$\hat t_8\,\hat t_8\,\hat R^4$ in 12d}
%----------------------------------------------------------------------

The $t_8$ tensor is defined as the product-of-metrics structure appearing
in the four-graviton amplitude (see \cite{Minasian:2015bxa}, appendix A):
\begin{equation}\label{eq:t8_def}
\begin{aligned}
    \hat t_8^{N_1\dots N_8}
    &= \tfrac{1}{16}\Bigl(
        -2\bigl(\hat g^{N_1N_3}\hat g^{N_2N_4}\hat g^{N_5N_7}\hat g^{N_6N_8}
        + \hat g^{N_1N_5}\hat g^{N_2N_6}\hat g^{N_3N_7}\hat g^{N_4N_8}
        + \hat g^{N_1N_7}\hat g^{N_2N_8}\hat g^{N_3N_5}\hat g^{N_4N_6}\bigr)\\
    &\quad +8\bigl(\hat g^{N_2N_3}\hat g^{N_4N_5}\hat g^{N_6N_7}\hat g^{N_8N_1}
        + \hat g^{N_2N_5}\hat g^{N_6N_3}\hat g^{N_4N_7}\hat g^{N_8N_1}
        + \hat g^{N_2N_5}\hat g^{N_6N_7}\hat g^{N_8N_3}\hat g^{N_4N_1}\bigr)\\
    &\quad -(N_1\leftrightarrow N_2) - (N_3\leftrightarrow N_4)
           - (N_5\leftrightarrow N_6) - (N_7\leftrightarrow N_8)
    \Bigr)\,.
\end{aligned}
\end{equation}

The 12d promotion of $t_8\,t_8\,R^4$ is
\begin{equation}\label{eq:t8t8R4_12d}
    S_{t_8}^{(3)}
    = \int f_0(\tau,\bar\tau)\;\hat t_8\,\hat t_8\,\hat R^4\;{*_{10}}\,1\,,
\end{equation}
where
\begin{equation}\label{eq:t8t8R4_contraction}
    \hat t_8\,\hat t_8\,\hat R^4
    = \hat t_8^{M_1\dots M_8}\,\hat t_{8\,N_1\dots N_8}\;
      \hat R^{N_1N_2}{}_{M_1M_2}\,
      \hat R^{N_3N_4}{}_{M_3M_4}\,
      \hat R^{N_5N_6}{}_{M_5M_6}\,
      \hat R^{N_7N_8}{}_{M_7M_8}\,.
\end{equation}

\paragraph{Goal.}
Expanding \eqref{eq:t8t8R4_contraction} using the 12d metric
decomposition \eqref{eq:12d_metric} yields the full 10d action
involving $R_{mnrs}$, $P_m$, $\bar P_m$, $D_m$, and $R_{mn}$.
This is equation (B.10) of \cite{Minasian:2015bxa}, which we aim to reproduce.
